\par\noindent\textbf{(1)} 稳态局部能量守恒 ∂\_xj\_q=0，Fourier定律 j\_q=−κ(T)T\_x，故 T\_x=−j\_q/κ(T)。从x=0积到x得 x/L=[∫\_\{T(x)\}\textasciicircum{}\{T\_h\}κ(s)ds]/[∫\_\{T\_c\}\textasciicircum{}\{T\_h\}κ(s)ds]，该积分隐式给出非线性温度剖面，要求κ(T)>0。
\par\noindent\textbf{(2)} 常热流 j\_q=(1/L)∫\_\{T\_c\}\textasciicircum{}\{T\_h\}κ(T)dT。若n≠−1，j\_q=κ₀T\_*\textasciicircum{}\{−n\}[T\_h\textasciicircum{}\{n+1\}−T\_c\textasciicircum{}\{n+1\}]/[(n+1)L]；n=−1则 j\_q=(κ₀T\_*/L)ln(T\_h/T\_c)。两式均为正，热从高温端流向低温端。
\par\noindent\textbf{(3)} 不可逆热传导局域熵产生密度 σ=κ(T)(T\_x/T)²≥0。对板积分并用j\_q常量、dx=−κ(T)dT/j\_q，得到 ∫₀\textasciicircum{}Lσ dx=j\_q(1/T\_c−1/T\_h)，恰等于两端热库熵流净增，与κ(T)具体幂次无关。
\par\noindent\textbf{(4)} 定义有效κ\_eff=Lj\_q/(T\_h−T\_c)，则κ\_eff是温度区间的算术积分平均，而非随便取κ((T\_h+T\_c)/2)。在ΔT→0时κ\_eff→κ(T\_0)，j\_q≈κ(T\_0)ΔT/L，回到常系数线性Fourier规律。
